\newpage
\thispagestyle{empty}
\vspace*{0.35\textheight}
\begin{center}
    {\Huge\textbf{CAPÍTULO I}} \\[0.5cm]
    {\Huge\textbf{MARCO INTRODUCTORIO}}
\end{center}

\newpage
\chapter{Marco Introductorio}

\section{Introducción}

La búsqueda de convocatorias públicas en el portal SICOES constituye una actividad relevante para empresas, consultoras y proveedores que participan en procesos de contratación estatal en Bolivia. Sin embargo, la exploración manual o basada únicamente en palabras clave presenta limitaciones cuando las necesidades del usuario no coinciden exactamente con la redacción administrativa del objeto de contratación.

\section{Planteamiento del problema}

Los mecanismos tradicionales de búsqueda se apoyan principalmente en coincidencias textuales. Este enfoque reduce la capacidad de recuperación cuando una convocatoria describe el mismo requerimiento con terminología distinta, sinónimos o expresiones técnicas sectoriales.

No obstante, en dominios administrativos como SICOES también existen consultas donde la coincidencia léxica exacta sigue siendo altamente efectiva. Por ello, el presente trabajo no asume a priori la superioridad universal de la búsqueda semántica pura, sino que compara de manera reproducible estrategias keyword, semánticas e híbridas para identificar la combinación más adecuada al dominio.

\section{Objetivo general}

Diseñar e implementar un sistema inteligente de recuperación semántica para convocatorias públicas del SICOES mediante embeddings, PostgreSQL con pgvector y Retrieval-Augmented Generation.

\section{Objetivos específicos}

\begin{enumerate}
    \item Extraer convocatorias públicas vigentes del portal SICOES.
    \item Limpiar y transformar los datos para construir un dataset orientado a recuperación semántica.
    \item Generar embeddings sobre las convocatorias.
    \item Implementar PostgreSQL con pgvector como base vectorial principal.
    \item Desarrollar una línea base de búsqueda tradicional usando SQL e \texttt{ILIKE}.
    \item Construir un prototipo de consulta tipo chatbot con LangChain.
    \item Evaluar búsqueda keyword, búsqueda semántica y una alternativa híbrida.
\end{enumerate}

\section{Justificación}

La propuesta es pertinente por su valor práctico y académico. En el plano aplicado, puede facilitar el acceso a oportunidades de contratación pública. En el plano metodológico, permite documentar un caso reproducible de integración entre procesamiento de lenguaje natural, bases vectoriales y generación aumentada por recuperación en un contexto boliviano.
